\section{Обзор литературы}

\subsection{Модели волатильности}

Исследование волатильности финансовых активов берёт начало с работ Энгла~\cite{Engle1982}
(модель ARCH) и Боллерслева~\cite{Bollerslev1986} (GARCH),
которые моделируют условную дисперсию доходностей как функцию прошлых наблюдений.
Однако GARCH-модели плохо работают на высокочастотных данных
и требуют сильных предположений о распределении остатков.

Ключевым вкладом в область прогнозирования волатильности стала модель HAR-RV
(Heterogeneous Autoregressive model of Realized Volatility), предложенная
Корси~\cite{Corsi2009}.
Модель основана на гипотезе рыночной гетерогенности: разные участники рынка
(краткосрочные трейдеры, институциональные инвесторы, пенсионные фонды) действуют
на различных временных горизонтах~--- дневном, недельном и месячном.
Спецификация HAR-RV:
\begin{equation}\label{eq:har}
    \mathrm{RV}_{t+1} = \beta_0 + \beta_d \, \mathrm{RV}_t^{(d)}
    + \beta_w \, \mathrm{RV}_t^{(w)}
    + \beta_m \, \mathrm{RV}_t^{(m)} + \varepsilon_{t+1},
\end{equation}
где $\mathrm{RV}_t^{(d)}$, $\mathrm{RV}_t^{(w)}$, $\mathrm{RV}_t^{(m)}$~---
среднеарифметические реализованные волатильности за последние 1, 5 и 22 торговых дня
соответственно.
Несмотря на линейность, модель~\eqref{eq:har} показывает хорошие
результаты на многих рынках \cite{Corsi2009}.

Для оценки качества прогнозов в работе используется QLIKE-функция
потерь~\cite{Patton2011}. Её удобство в том, что она остаётся корректной,
даже когда вместо ненаблюдаемой «истинной» волатильности подставляется
её оценка (прокси):
\begin{equation}\label{eq:qlike}
    \mathrm{QLIKE} = \mathbb{E}\!\left[
        \frac{\mathrm{RV}_t}{\hat{\sigma}_t^2}
        - \ln\!\frac{\mathrm{RV}_t}{\hat{\sigma}_t^2} - 1
    \right].
\end{equation}

\subsection{Attention и sentiment в прогнозировании}

Тетлок~\cite{Tetlock2007} одним из первых продемонстрировал, что тональность
публикаций в Wall Street Journal предсказывает доходности и объёмы торгов на фондовом рынке.
Да, Энгельберг и Гао~\cite{Da2011} предложили использовать индекс поисковой активности Google
как прямую меру внимания розничных инвесторов и показали его связь с будущими доходностями.

Работа Halousková и Lyócsa~\cite{Halouskova2025} лежит в основе настоящей работы.
Авторы строят вектор «макроэкономического внимания»~--- набор признаков,
отражающих интерес инвесторов к макроэкономическим анонсам (инфляция, ставка ФРС,
ВВП и др.)~--- и показывают, что его добавление в HAR-модель заметно улучшает
прогноз волатильности акций США.
Настоящая работа адаптирует этот подход к российскому рынку драгоценных металлов,
заменяя американские анонсы на решения ЦБ РФ, геополитические события
и данные русскоязычного новостного потока.

Для извлечения sentiment из финансовых текстов широко применяются
модели на основе BERT~\cite{Araci2019}.
FinBERT, обученный на финансовых новостях и аналитике,
работает заметно лучше словарных методов
при классификации тональности финансовых текстов.
В настоящей работе используется русскоязычная модель ruBERT
для обработки новостей российских СМИ.

\subsection{Машинное обучение для прогнозирования волатильности}

Ченом и Гестрином~\cite{Chen2016} была предложена библиотека XGBoost~---
реализация градиентного бустинга с регуляризацией,~--- ставшая стандартом
в задачах прогнозирования на табличных данных.
Применительно к финансовым временным рядам XGBoost позволяет нелинейно
комбинировать HAR-признаки с sentiment-вектором без явного задания функциональной формы.

Рекуррентные нейронные сети, в частности LSTM~\cite{Hochreiter1997},
способны улавливать долгосрочные зависимости в последовательностях,
что делает их естественным выбором для моделирования волатильности
с явной последовательной структурой.

\subsection{Волатильность рынка драгоценных металлов}

Рынок драгоценных металлов характеризуется рядом особенностей,
отличающих его от рынков акций:
высокой чувствительностью к геополитическим событиям и инфляционным ожиданиям,
ролью золота как «safe haven» актива и значительным влиянием
курса доллара США на рублёвые цены инструментов MOEX.
Российский рынок металлов после 2022 года приобрёл дополнительную специфику
в связи с ограничениями на экспорт и реструктуризацией торговых потоков,
что создаёт нетривиальные задачи для прогнозирования.
